\needspace{8\baselineskip}
\section{Discussion}
\label{sec:discussion}

\chg{This study examined how far a 16-electrode forearm sEMG array could be reduced while preserving subject-independent gesture-recognition performance under a prespecified engineering tolerance. The principal finding is deliberately modest: four electrodes could be removed, whereas five could not. The smallest non-inferior configuration therefore contained \KStar\ electrodes, representing a \KStarReductionPct\% reduction from the complete array. Importantly, this result was not driven by a single irregular point on the performance curve. All configurations from 12 to 15 electrodes remained within the non-inferiority margin, and the minimum acceptable count was unchanged after correction for multiple testing. At the same time, every reduction produced some numerical loss in Macro-F1, and halving the array incurred a substantially larger decrease. The full montage therefore contains some redundancy, but considerably less than might be inferred from user-dependent electrode-selection studies.}

\chg{This distinction between user-dependent and subject-independent reduction is central to the interpretation of the result. A sparse configuration optimised for a particular individual can exploit that person's anatomy, electrode placement and activation patterns. A layout intended for previously unseen users must instead preserve information across inter-individual variability. The present result therefore should not be interpreted as contradicting studies reporting successful control with substantially fewer electrodes \citep{phinyomark2018emg,qu2021reduce}; rather, it addresses a different design question. Under the current montage, gesture set and cross-user protocol, a moderate reduction was supported, whereas aggressive sparsification was not.}

\chg{Cross-fitted attribution provided useful information for selecting which electrodes to retain, but its benefit was more limited than a comparison with random subsets alone might suggest. Attribution-selected configurations outperformed the mean of random configurations at every tested subset size, and deletion with retraining showed that removing highly ranked electrodes was more damaging than removing weakly ranked ones. Together, these findings support the functional relevance of the ranking. However, the attribution-based configuration was rarely among the very best random subsets for an individual participant, and exploratory comparisons showed no statistically detectable advantage over permutation importance or an evenly spaced geometric layout at most tested subset sizes. The clearest separation occurred at six electrodes, where the attribution-guided ranking scored higher than both the geometric layout and permutation importance in uncorrected exploratory contrasts.}

\chg{This pattern suggests that two effects may be operating simultaneously. The forearm montage appears to contain broad spatial redundancy, making circumferential coverage an effective baseline when a relatively large number of electrodes is retained. At six electrodes, the identity of individual recording locations appeared to matter more. No advantage over the geometric layout was detected at four electrodes, however, so the benefit of attribution did not simply increase as the array was reduced. This interpretation remains exploratory: the study did not formally test equivalence between attribution-based and geometric layouts, nor did it test whether the geometric layout itself satisfies the non-inferiority criterion. The absence of a statistically detectable difference should therefore not be interpreted as evidence that the two approaches are interchangeable.}

\chg{A second important result is negative. Penalising attribution instability did not improve reduced-array prediction. The cross-fitted attribution profiles were already highly similar across inner participant groups, leaving little instability for the regularisation term to penalise; the penalty changed the top-12 set in at most one of 43 folds. Consequently, the stability-adjusted rankings were almost identical to the rankings based on attribution magnitude alone. The post-hoc participant-level extension made between-person variability more visible but reached the same practical conclusion: a more sensitive estimate of attribution heterogeneity did not translate into improved electrode selection. Within the present dataset and sample size, the tested stability formulation therefore added no detectable ranking information beyond plain cross-fitted TreeSHAP magnitude, which is the simpler option. This result does not exclude a role for stability regularisation \citep{xiang2026roshap} in smaller cohorts, more heterogeneous populations, or settings in which attribution estimates vary more strongly across subjects.}

\chg{The spatial consistency of the ranking also warrants careful interpretation. Several electrode positions were repeatedly ranked near the top across outer folds and individual participants, indicating reproducible regional differences in discriminative information. However, these positions cannot be interpreted directly as specific muscle contributions. GRABMyo documents the electrode montage more reliably than the underlying muscle boundaries, no imaging-based electrode-to-muscle registration was available, and each surface recording reflects volume-conducted activity from multiple sources \citep{farina2004extraction}. The present results therefore support a positional interpretation (some circumferential regions are consistently more informative than others) but not muscle-level localisation.}

\chg{The exploratory analyses further define the scope of the primary result. NinaPro DB5 did not provide a decisive external replication because its ten participants yielded wide uncertainty and the candidate non-inferior configuration did not survive multiplicity correction. Similarly, the cross-session analysis did not reveal an additional performance penalty on later recording days, but this should not be interpreted as evidence that electrode displacement or repeated donning is inconsequential \citep{hargrove2006effect,kaufmann2010fluctuating,vidovic2016improving}. All conditions remained subject-independent, and training on multiple sessions also increased the amount of available training data. Most importantly, the bipolar re-referencing analysis did not reproduce the degree of reduction observed for monopolar electrodes. Removing even one of eight derived bipolar pairs failed the primary non-inferiority criterion. The headline \KStarReductionPct\% reduction is therefore specific to monopolar recording sites on the analysed two-ring montage and should not be translated directly into a claim about reducing physical bipolar sensors.}

\chg{Several limitations constrain the generality of these findings. All participants were able-bodied, and the analysis was entirely offline. Zero-phase filtering is non-causal, and the present evaluation therefore does not establish real-time control performance. Effects of fatigue, posture, electrode displacement and repeated donning were not isolated. The electrode rankings were derived from models trained on the complete array and may differ from rankings learned directly under sparse acquisition. Computational constraints also required reduced training samples for some classifier and selection-method comparisons.} The support vector machine result is therefore likely a lower bound, and the H2 and H4 comparisons used a lighter model than the performance curve. Training-sample sizes were chosen on two folds using held-out outcomes, and $\lambda$ was not selected. \chg{In addition, several secondary test families and exploratory analyses were specified after the primary results were available, as documented in supplementary methods~S1, and should be interpreted accordingly;} the fixed-layout result (supplementary methods~S4) is a single-split consistency check only, and the window-length and 17-class analyses listed in the plan were not run. \chg{Finally, the independent replication dataset was small, and the bipolar analysis used signals derived mathematically from monopolar recordings rather than measurements from a physical bipolar montage.}

\chg{Taken together, the findings support a restrained design conclusion rather than a claim for aggressive electrode minimisation. For subject-independent recognition, the relevant question is not whether some sparse configuration can perform well for particular users, but whether the performance loss associated with a common reduced configuration remains within a tolerance defined before observing the result. Under that criterion, this study supports a moderate reduction from 16 to 12 monopolar electrodes. Cross-fitted attribution offers a leakage-controlled way to identify informative sites, but the results also show that its practical advantage depends on how strongly the array is reduced and may be modest when broad spatial coverage is already preserved.}
